\chapter{INTELLIGENCE}\label{intelligence}

\hfill\break

Un paio di giorni dopo, Lester Cornhut mi fermò durante un
pattugliamento mattutino per dirmi eccitato che il governatore regionale
sarebbe stato a casa sua alle 16:00 in punto. Si tamburellò con le dita
il polso come se avesse indosso un orologio e mi chiese se potessi
andare nella sua umile dimora. A volte, il traduttore faceva le cose in
grande: sono certo che non avesse davvero detto ``umile dimora'' in
ruhar. Ero tentato di lasciare i criceti ad aspettarmi fino alle calende
greche, ma mia madre mi ha insegnato a essere educato e i militari mi
hanno messo in testa l'abitudine di essere puntuale, per tutto. Alle
16:00 in punto ora locale arrivai alla residenza di Cornhut per
incontrami con il governatore regionale, Lester Cornhut, la signora
Cornhut e i piccoli Cornhut. Baker aspettò fuori, Sanchez e Chen erano
alla postazione di comando con le armi pesanti, per ogni evenienza. La
mia visita era stata autorizzata, anche incoraggiata, dal quartier
generale del plotone. Lasciai fucile, casco, giubbotto antiproiettile e
occhiali alla postazione di comando, portando solo il mio zPhone e
l'arma da fianco. L'arma l'avevo perché mi sarei sentito nudo senza, e
perché avrebbe ricordato, ai criceti e a me, che l'Unef era una forza di
occupazione. Enfasi su ``forza'', se necessario.

Ah, e avevo portato con me una bustina di tè. E un pacchetto di
zucchero, nel caso ne avessi avuto bisogno. Lester mi salutò entusiasta,
la mia ipotesi era che il governatore regionale non venisse spesso in
visita in una città di campagnoli come Teskor, è probabile che da quel
posto passassero più tornado che governatori.

C'erano soltanto tre Cornhut in famiglia: Lester, sua moglie e un
figlio, che riconobbi dalla scuola. «Grazie per il pallone», disse il
bambino in inglese, o almeno così sembrava. Fui davvero colpito dal
fatto che il piccoletto peloso si fosse sforzato di memorizzare quelle
parole, perciò dissi: «Ta», e feci un breve inchino. Lui sorrise, poi
Lester lo cacciò via e versò acqua calda in due tazze su una sorta di
tavolino da caffè davanti al divano. Infine, fece un profondo inchino e
se ne andò, lasciandomi solo con il governatore regionale. Assomigliava
a qualsiasi altra ruhar di sesso femminile, anche se era chiaro che non
era una contadina; indossava una camicetta e una gonna lunga, di un
materiale simile alla seta, orecchini e collana all'apparenza costosi,
non che avessi idea di come distinguere un gioiello da un altro. Era
seduta su un divano verde scuro, con sopra un tappetino fiorato; io mi
sedetti su una sedia all'estremità del tavolino, di sicuro non mi sarei
seduto su quel cazzo di divano accanto a lei. Da una scatola sul tavolo
prese lentamente col cucchiaio quelle che sembravano foglie di tè sfuse,
le mise in una cosa a forma di uovo con un sacco di buchi e una
catenella sottile, che calò con delicatezza nella tazza da tè. Mia nonna
aveva un coso ovale da tè, o come diavolo si chiama, quasi identico.

«Sono il sergente Joe Bishop», dissi e misi la mia bustina di tè in una
tazza. Il mio modo di fare il tè era molto meno elegante del suo. Questo
m'irritò, per una qualche ragione. Conoscevo il motivo per cui aveva
fatto il tè così lentamente: di proposito, era per prendersi il tempo di
studiarmi.

Sembrava divertita: «Sono il governatore regionale di Lesscorta, mi
chiamo Bahturnah Lohgellia». Fece una pausa, in attesa che il traduttore
si rimettesse al passo. «La sua gente che occupa la mia città mi
chiama», si fermò di nuovo, «la Bürgermeister.» Le brillavano gli occhi,
era chiaro che si aspettava che lo trovassi divertente. «Penso che
``Bürgermeister'' sia più facile da pronunciare del mio nome, no?»

Ero divertito. «Ta.» Quali che fossero le truppe occupanti la sua città,
dovevano essere state di stanza in Germania a un certo punto. Il
Bürgermeister è una specie di sindaco in terra tedesca. Almeno credo.
Voglio dire, la cosa più simile all'essere stato in Germania che abbia
mai fatto è stata sedere in un aeroplano da trasporto mentre faceva
rifornimento di carburante sulla pista dell'aeroporto di Rhine-Main.

La Bürgermeister mi fece cenno di togliermi microfono e auricolare. Io
mi indicai la bocca, poi l'orecchio. «Ne ho bisogno perché possiamo
parlare.»

Con mia grande sorpresa, lei rispose in un inglese molto lento, cauto e
stridulo: «Sergente Joe Bishop, per favore metta via quello. Può usare
questo». In mano aveva quello che sembrava uno zPhone kristang, ma un
po' più piccolo. Me l'offrì, con il suo microfono e il suo auricolare.
«Le spiegherò», disse lentamente, come se avesse memorizzato solo poche
frasi in inglese e trovasse le nostre parole difficili da pronunciare.

Esitai, poi annuii. I regolamenti dell'esercito dicevano che non dovevo
separarmi dal mio zPhone mentre ero in servizio, ma questa era
un'opportunità per mettere le mani sulla tecnologia ruhar e, forse, per
ottenere qualche informazione sui criceti. Era soprattutto per questo
che l'Unef aveva creato le squadre integrate di osservazione. Uscii,
diedi il mio zPhone a Baker, tornai e usai il dispositivo ruhar.

«Grazie», disse la Bürgermeister in inglese, poi si mise l'auricolare e
tornò a usare il traduttore. «È comodo?» Anche la versione
computerizzata della sua voce aveva un tono stridulo.

Guardai lo schermo del dispositivo aspettandomi che mi chiedesse di
pronunciare ad alta voce un paio di pagine di testo, in modo da poter
acquisire il mio particolare schema vocale, ma visualizzavo solo la
scritta rossa ``Pronto'' in grassetto. Il che mi rese sospettoso.
«Questo come fa a sapere come parlo?»

La Bürgermeister sorrise. «Il suo modo di parlare è stato registrato,
analizzato, tradotto e programmato nel dispositivo che ha in mano. Non
volevamo che perdesse tempo a configurare il dispositivo, perché abbiamo
molte cose di cui parlare.»

Quindi mi stavano spiando. Giusto, ovvio che volessero informazioni sul
loro avversario. «Perché non posso usare il mio\ldots», stavo per dire
zPhone, ma pensai che non si traducesse bene in ruhar, «la mia radio
tattica?»

Di nuovo sorrise. Un sorriso amichevole, ma anche un sorriso alla ``So
qualcosa che tu non sai''. Stava diventando un po' irritante. «I
kristang non vi hanno fornito armi avanzate», indicò la mia arma da
fianco, «non è curioso di sapere come mai vi abbiano invece fornito
attrezzature di comunicazione avanzate, utilizzate da ogni essere umano
su questo pianeta?»

Questo mi fece riflettere. «Ah», cazzo, non ero un esperto di
comunicazioni. «È probabile che sia perché, ehm, ogni paese sulla Terra
usa apparecchi radio e computer diversi e, dato che non abbiamo i nostri
satelliti di comunicazione qui», alzai gli occhi al soffitto, «l'unico
modo in cui possiamo parlare tra di noi è con radio comuni.» Come lo
dissi, aveva senso, anche se prima non ci avevo pensato molto su.

Niente sorriso stavolta, la Bürgermeister scosse la testa da un lato
all'altro. Mi affascinava il fatto che, tra due specie aliene, il
linguaggio del corpo sembrasse essere universale. Non c'era bisogno di
traduzioni in quel caso. «Forse questa è una scusa conveniente. La vera
ragione è che tutte le vostre comunicazioni passano attraverso una rete
kristang. I kristang controllano ogni parola che dite, catturano tutte
le vostre trasmissioni di dati, tracciano ogni vostro movimento, ovunque
sul pianeta. Possono anche bloccare tutte le vostre comunicazioni,
quando vogliono. Sergente Bishop, quando uscirà di qui, riferirà la
nostra conversazione ai suoi servizi segreti militari. Le suggerisco di
condurre questa conversazione di persona, piuttosto che via radio. A
meno che non voglia che le lucertole la ascoltino.»

Era tutto vero, credo. Noi usavamo solo attrezzatura kristang per
comunicare. Ma, perché no? Era gratuita, funzionava alla grande, ci
risparmiava lo sforzo d'installare ponti radio in tutto il pianeta e di
provare a far parlare le radio americane con le radio cinesi e indiane.
Inoltre, i kristang erano nostri alleati. «Forse ha ragione. E allora? I
kristang sono nostri alleati.»

Quel sorriso irritante tornò. «Alleati? Le alleanze si stringono tra
pari. I kristang sono i vostri patroni e voi siete la loro specie
cliente. I loro animali domestici. O schiavi.»

«E voi siete i nostri nemici», dissi senza pensarci. Mia madre non ne
sarebbe stata felice: ero ospite a casa dei criceti a bere il tè e non
mi stavo comportando in modo gentile.

«Non abbiamo nessuna intenzione di essere vostri nemici. Per favore, mi
lasci parlare», lei alzò una mano. «Posso capire se è arrabbiato, c'è
molto che deve sapere, molto che i kristang non le hanno detto, o su cui
hanno mentito. La mia gente conosce la vostra specie da più di mille
anni. Le nostre sonde ad ampio raggio vi hanno trovato e hanno collocato
satelliti stealth in orbita attorno al vostro pianeta per osservarvi. Mi
dica, non si è chiesto perché, visto che questa guerra va avanti da
migliaia di anni, la Terra non sia stata attaccata già molto tempo fa?»

La mia mascella si schiuse. Aveva ragione, quel punto non mi tornava,
non mi tornava per niente, e non ero l'unico ad aver posto quella
domanda. Perché aspettare di attaccare ora che gli umani avevano le
testate nucleari? Perché non attaccare molto tempo prima, quando gli
vivevano nelle caverne e pensavano che il fuoco fosse l'apice della
tecnologia? Non aveva alcun senso. I ruhar non erano stupidi, quindi?
Perché? Se al quartier generale dell'Unef lo sapevano, non me l'avevano
detto.

Poiché la necessità di usare un traduttore faceva sì che la
Bürgermeister parlasse molto lentamente, farò un riassunto di quello che
disse. Il motivo per cui la Terra non era stata attaccata in precedenza
è semplice. Fino a poco tempo prima, si trovava fuori dal range pratico
di volo, per entrambe le fazioni in guerra. Era possibile, per astronavi
specializzate ad ampio raggio, raggiungere la nostra misera parte del
Braccio di Orione, solo che la Terra non valeva il tempo o le spese. Per
qualche ragione, persino i maxolhx e i rindhalu non sapevano che, di
tanto in tanto, i wormhole si spostavano. Accadeva a random o a cascata
nell'intero quadrante della galassia. Uno spostamento poteva avvenire in
un paio di decenni, o richiedere centinaia o migliaia di anni. Un
wormhole A, un tempo connesso a un wormhole B, a un tratto si connetteva
invece a un wormhole C, e B si spostava per connettersi a D.

Oppure un wormhole smetteva del tutto di funzionare, mentre uno prima
sconosciuto poteva attivarsi all'improvviso. Un pianeta che era
importante dal punto di vista strategico poteva diventare una zona
isolata, dopo lo spostamento di un wormhole. All'inverso, un pianeta che
era troppo lontano da un wormhole perché qualcuno se ne curasse, poteva
diventare all'improvviso un'importante base di sosta per entrambe le
fazioni in guerra. È quello che era accaduto alla Terra. Vivevamo
felicemente da soli nell'entroterra galattico, distanti da qualsiasi
wormhole e all'estremità del territorio dei ruhar, ma abbastanza lontani
da far sì che i ruhar non ritenessero conveniente percorrere la distanza
enorme che li separava dalla Terra. Poi c'era stato uno spostamento e un
wormhole da tempo addormentato si era attivato. Un'estremità del
wormhole era poco distante dalla Terra, l'altra si trovava nel
territorio dei kristang. Quel wormhole era abbastanza vicino alla Terra
perché una nave madre thuranin potesse raggiungerla in una settimana
circa, il che aveva fatto della Terra un buon posto agli occhi dei
kristang, che vi avrebbero ottenuto un punto d'appoggio nel territorio
dei ruhar. Prima dello spostamento del wormhole, la Terra era troppo
lontana per una campagna militare dei ruhar e impossibile da raggiungere
per i kristang.

«Aspetti un minuto», dissi mentre digerivo quello che la Bürgermeister
mi aveva detto. «Avete attaccato la Terra. Se siamo troppo lontani per
un'effettiva campagna militare, perché ve ne siete presi la briga? Deve
essere stato uno sforzo enorme.» Pensavo alla complicata logistica della
nostra occupazione di Paradiso, e noi avevamo un wormhole che copriva la
maggior parte del percorso degli approvvigionamenti. Se i ruhar dovevano
scarpinare per tutta la strada, come direbbe mio nonno, come cazzo
avevano fatto? E perché attaccare un pianeta, se non si può sostenere
una forza lì? A meno che noi non fossimo una minaccia per loro.

La Bürgermeister annuì. Era chiaro che era preparata a questa domanda.
«Non sono sicura che il suo traduttore stia raccogliendo le mie parole
con accuratezza, perciò, per favore, mi avvisi se non capisce. Quando ci
siamo resi conto che lo spostamento del wormhole aveva aperto il vostro
pianeta allo sfruttamento da parte dei kristang, abbiamo deciso di agire
per primi e i nostri patroni, gli jeraptha, hanno radunato le navi ad
ampio raggio che erano disponibili a rispondere. Una volta saputo per
certo che i kristang intendevano occupare il vostro pianeta, abbiamo
danneggiato la vostra infrastruttura industriale per fare della Terra
una base di sosta meno appetibile agli occhi dei kristang. È stato un
attacco rapido da parte di un piccolo numero di navi, non potevamo
lanciare un attacco su larga scala a quella distanza. Le nostre navi
hanno viaggiato per cinque dei vostri mesi per raggiungere la Terra e
cinque per tornare.»

Pensai che era una grande stronzata. «E siete arrivati nello stesso
istante dei kristang? Bene.» Poi mi resi conto che ``bene'' poteva
essere tradotto come ``esatto''. Ci era stato detto di parlare in un
linguaggio semplice coi ruhar, perché idiomi, slang e linguaggio
emozionale come il sarcasmo potevano perdere senso attraverso la
barriera della traduzione tra le specie. Sebbene basato sulle mie
interazioni con i ruhar, il linguaggio del corpo era universale, almeno
tra le specie bipedi. «Quello che volevo dire è che dubito che ciò che
ha detto sia preciso.»

«Non siamo arrivati nello stesso istante dei kristang», spiegò lei con
quello che interpretai come un sorriso di lieve condiscendenza. «Le
nostre navi sorvegliavano il wormhole, in caso i thuranin vi avessero
inviato una forza, sapevamo che stavano sondando la zona, ma non eravamo
del tutto sicuri delle loro intenzioni. La Terra non è l'unico pianeta
abitabile vicino al nuovo wormhole. Dopo che la forza d'invasione vi era
entrata, abbiamo aspettato per tre settimane che le navi thuranin si
riunissero e completassero il viaggio. Quando abbiamo capito che il loro
obiettivo era la Terra, le nostre navi pesanti li hanno ingaggiati,
mentre le nostre forze di attacco hanno lanciato colpi di precisione
sulla Terra. Non abbiamo colpito le città, il nostro scopo non era
uccidere la sua gente, era ridurre la vostra capacità industriale, così
i kristang non avrebbero avuto vita facile nel mantenere un punto
d'appoggio ai margini del nostro territorio.»

«Avete ucciso un sacco di persone», gridai con rabbia in risposta.
«Umani.»

«Da soldato, ha familiarità con l'espressione ``danno collaterale''?»
Aspettò che annuissi. Annuire era uno di quei linguaggi del corpo
universali. «I nostri obiettivi erano le centrali elettriche, le
fabbriche che producono determinati tipi di attrezzature e gli impianti
industriali che creano materiali come i metalli. Tutte infrastrutture
che avrebbero potuto essere utilizzate dai kristang. Quando abbiamo
attaccato le vostre centrali a fissione», pensai che il traduttore
intendesse ``a energia nucleare'', «abbiamo colpito il centro di
distribuzione elettrica nei pressi del reattore. Siamo stati attenti a
non colpire il reattore nucleare. Non volevamo contaminare il vostro
pianeta e causare morti e sofferenze senza motivo. I soldati ruhar che
sono stati costretti ad atterrare nel suo villaggio stavano per
distruggere il centro di distribuzione elettrica di una centrale a
fissione in un posto chiamato Connecticut.» Poiché non c'era un
equivalente ruhar di Connecticut, la parola mi giunse nella sua voce da
roditore un po' stridula, senza traduzione. «Non potevamo rischiare di
colpire quella centrale dall'orbita, neppure con missili intelligenti.»

Questo mi fece riflettere. Diceva la verità sugli attacchi alle centrali
nucleari, nessuno dei reattori in tutto il pianeta era stato colpito.
«Non volevate contaminare un pianeta che volevate occupare!» Feci un
cenno con la mano per zittirla, mentre cercava di parlare: «No, ho
chiuso con le sue bugie». Sbattei la tazza da tè con rabbia sul tavolo,
mi alzai e feci un rigido inchino. «Grazie, signora.» Quando uscii, vidi
un gruppo di bambini criceti giocare con il pallone da calcio che avevo
portato, e non riuscivo a decidere se ero più arrabbiato con i criceti
per avere cercato di giustificare l'attacco alla Terra, o per essermi
concesso di essere morbido con il nemico. Uno dei bambini criceto mi
salutò con la mano e io lo ignorai. Più tardi mi sarei sentito una merda
per questo, erano bambini alieni, ma erano pur sempre bambini. Non
avevano attaccato la Terra. Volevo solo qualcuno con cui prendermela.

\hfill\break

Restai incazzato per quasi una settimana, arrabbiato per essermi
lasciato manipolare dalla Bürgermeister. Arrabbiato perché aveva cercato
di usarmi per fornirci informazioni fallaci, per seminare dissenso tra
le nostre fila, per condurre una campagna di manovra psicologica contro
di noi. Non avrebbe funzionato. Quando alla fine ne discussi con un
ufficiale, durante una delle visite del tenente Charles, gli parlai
faccia a faccia e feci il gesto di sottrargli lo zPhone: con mia grande
sorpresa lui annuì in silenzio e lasciò il suo zPhone fuori dalla tenda.
«Cosa c'è, sergente?»

«Signore, quel governatore regionale criceto mi ha dato informazioni su
come funzionano i wormhole. Potrebbero essere tutte stronzate, ma ho
pensato di dover fare rapporto ai superiori. E ha detto che non dovrei
parlarne allo zPhone, perché i kristang potrebbero essere in ascolto.»

Lui annuì. «Il comando si è preoccupato di questo. I nostri alleati
controllano tutte le nostre comunicazioni su Paradiso. E il nostro
trasporto. E le nostre scorte di cibo e tutto il resto. Va bene, cosa
sai dei wormhole?»

Glielo dissi, come meglio ricordavo. Lui annuì. «Hai ragione, è
probabile che siano stronzate, ma riferirò al quartier generale del
battaglione.»

\hfill\break

Lester Cornhut mi aveva detto che la Bürgermeister voleva parlare di
nuovo con me, ed era stato insistente, quasi implorante. Mi dispiaceva
per il criceto, ma ero ancora incazzato e non avevo nessuna intenzione
di farmi di nuovo manipolare da una donnola bugiarda, o da un criceto.
Secondo Lester Cornhut, il governatore regionale aveva fatto un viaggio
speciale a Teskor solo per parlare con me. Gli dissi grazie, ma no.
Sembrava ferito.

\hfill\break

Tre giorni dopo, un paio di Crivee si avvicinarono alla nostra
postazione di comando in una nuvola di polvere. Io ero nel nostro campo
da basket per una partita a horse¹⁷ con Baker quando arrivarono;
entrambi ci dannammo per metterci le magliette in fretta. Un maggiore
saltò fuori dal Crivee principale, accompagnato dalla sicurezza.
«Sergente Bishop?»

Feci il saluto. Somigliava vagamente all'ufficiale dei servizi segreti
della nostra brigata, di cui avevo visto soltanto delle fotografie. «Sì,
signora.»

«Sono il maggiore Perkins. Ho bisogno di parlare con te, dentro.» Fece
il gesto, ormai familiare, di togliermi lo zPhone.

Una volta all'interno, fui felice che avessimo dato una rinfrescata al
posto quella mattina, ci stavamo abituando a portare fuori la nostra
spazzatura in mezzo al nulla. «Bishop, le tue informazioni sui wormhole
hanno sollevato un polverone al quartier generale dell'Unef. Qualcuno là
ha riferito l'informazione sullo spostamento dei wormhole ai kristang e
loro hanno perso le staffe, volevano sapere dove avessimo sentito quelle
menzogne offensive e hanno negato tutto. Il che significa che è vero e
le maledette lucertole stanno mentendo in proposito.» Era la prima volta
che sentivo un superiore parlare contro i nostri alleati, mi sorprese.

«Lo hai sentito da una ruhar? Voglio incontrarla.»

«Ehm, signora, è complicato. Lei non vive qui al villaggio, è una specie
di governatore regionale, o qualcosa del genere. Passa di qui una volta
a settimana e ci siamo incontrati a casa del suo amico, per il tè.» Mi
sentivo un idiota a dire che mi ero seduto a bere il tè con il nemico.
«È una cosa da criceti, loro prendono il tè quando s'incontrano.»

«Va bene, quando tornerà?»

«Dopodomani, si presume, ma ho detto al suo amico qui che non volevo
parlare di nuovo con lei. Posso far sapere che ho cambiato idea?»

«Cazzo, Bishop, non funzionerà, devo incontrarla. Tu non sai che domande
abbiamo bisogno di farle. L'Unef ha molte cose di cui vuole parlare.»

«Non le faccio domande. Lei parla e io ascolto. Lei parla di qualunque
cosa voglia. Ho pensato che mi stesse prendendo per il culo, tipo una
manovra psicologica nello stile dei criceti.» Cazzo, la Bürgermeister
aveva detto la verità?! «Vuole sapere tutto quello che mi ha detto?»

«Sì.» Tirò fuori tablet e microfono dallo zaino. «Registrerò quello che
dici con questo. Con questo e non con uno zPhone, ti è chiaro?»

«Sì, signora. Forte e chiaro.» Soprattutto da quando la Bürgermeister mi
aveva detto di non fidarmi delle apparecchiature di comunicazione
fornite dai kristang.

Posò il tablet su un tavolo, collegò il microfono e si guardò attorno
nella nostra piccola postazione di comando. Era un bene che avessimo
riordinato e pulito il posto, cioè strofinato da cima a fondo, non solo
spazzato. Avevamo anche spostato i mobili e usato una specie di
lucidatrice per pavimenti che i criceti avevano lasciato in un armadio.
La lucidatrice aveva funzionato alla grande dopo che ci avevamo giocato
per un'ora, perché i criceti non avevano lasciato istruzioni utili. Il
maggiore Perkins annuì d'approvazione: «Ti piace qui, sergente? Sono un
po' isolate, queste squadre integrate di osservazione».

«Mi piace, signora. Sono un sergente novellino, perciò essere in mezzo
al nulla significa che posso fare errori senza che un sottotenente mi
stia col fiato sul collo.» Mentre parlavo, mi venne in mente che il
maggiore Perkins stesso doveva essere stata sottotenente un tempo. «Ehm,
non stavo\ldots»

«Rilassati, sergente.» Rise. «Quando ero sottotenente, ero forse
l'ufficiale più stupido nella storia dell'esercito degli Stati Uniti, e
questo la dice lunga. Certo, allora non lo sapevo. Questo posto può
anche essere in culo a Nettuno, ma goditelo finché puoi, a un sacco di
capisquadra di gruppi di fuoco piacerebbe avere questa opportunità di
uscire da soli. E se le informazioni che stai fornendo sono solide come
pensa l'Unef, stai andando alla grande. Cominciamo», aprì un'app per la
registrazione vocale sul suo tablet.

\hfill\break

Il maggiore Perkins tornò due giorni dopo, ma la Bürgermeister no.
Perciò, ce ne stavamo seduti a pettinare giaguari per la maggior parte
della giornata, finché un bambino criceto, in sella a una bicicletta
elettrica, non venne a cercarmi per trasmettermi il messaggio che la
Bürgermeister era occupata e avrebbe voluto incontrami il giorno
successivo. Me e soltanto me. Il maggiore Perkins non ne fu contenta, ma
capì l'antifona e mi diede una lista di domande che dovevo imparare a
memoria.

Fu così che iniziò la mia breve carriera da ufficiale dell'intelligence.
La lista di domande dell'Unef non portò a niente di buono; sembrava che
la Bürgermeister avesse un programma con quello di cui voleva parlare
ogni volta che ci incontravamo, e si stava attenendo con rigore a
quello. L'avrei incontrata due volte a settimana, e il maggiore Perkins,
talvolta accompagnata da un capitano dalla divisione del quartier
generale, arrivava il giorno successivo per interrogarmi. Mi sembrava di
essere un bambino che passava compiti al liceo, ma cosa avrei potuto
fare?

Durante il nostro quarto incontro, dopo esserci sbarazzati dell'intera
cerimonia del tè, posi una domanda alla Bürgermeister. «Perché io?»,
chiesi. «Perché non parla con un ufficiale dell'intelligence di una
delle nostre forze? Se non vuole parlare con un americano, può parlare
con inglesi, francesi, indiani o cinesi.»

La Bürgermeister si scurì in volto: «Non voglio essere interrogata. La
mia intenzione è fornire informazioni, per correggere le bugie che i
kristang le hanno detto e spiegare cose di cui i kristang non vogliono
parlare. Non parlerò con nessun altro a parte lei».

«Non ha risposto alla mia domanda. Perché io? Sono l'unico sergente di
grado inferiore. Il villaggio è a grande distanza da dove vive lei,
perché non parla con qualcuno là?»

Di nuovo, ebbi la sensazione che fosse preparata a ogni mia singola
domanda. La Bürgermeister era un criceto molto intelligente, su questo
non c'erano dubbi. Fece un altro sorriso, non avrei potuto dire se fosse
genuino, oppure quello falso tipico dei politici. «Sono venuta qui
perché volevo incontrare lei, Joe Bishop, quello che a volte chiamano
Barney.» Il traduttore lasciò che il mio nome e ``Bar-ney'' mi
giungessero nella sua voce stridula. Inclinò la testa, altra espressione
universale del linguaggio del corpo tra i bipedi. Cazzo, la sua fonte
d'informazioni era buona. «Ero curiosa di sapere che tipo di umano
avesse ferito un soldato ruhar e ne avesse catturato un altro.»

«Sono stato aiutato.»

«Aiutato da civili e, stando ai rapporti, non aveva armi militari.»

«Non sempre i rapporti sono accurati.» Non volevo fornirle informazioni
sulle tattiche che avevo utilizzato allora, o qualunque altro dettaglio
potesse risultare utile al mio nemico. Lei era là per fornirmi
informazioni, si presume, non per ascoltarmi.

«Questi rapporti lo sono. Un altro motivo per cui sto parlando con lei è
il rapporto del soldato ruhar che ha catturato.»

Non potei nascondere la sorpresa sulla mia faccia.

Lei continuò con un altro sorriso. «Non lo sa? È stato rilasciato, in
cambio di un soldato kristang che avevamo fatto prigioniero. Il nostro
soldato ha riferito di essere stato catturato da civili umani, cosa che
un soldato non ammette volentieri. E ha detto che lo avete trattato
bene. Lo apprezziamo.»

«Meglio che potevo. Non avevamo cibo che potesse mangiare.» Mi ero
chiesto cosa ne fosse stato del criceto dopo che la guardia nazionale lo
aveva portato via, mi ero fatto l'idea che lo avessimo consegnato ai
kristang prima che morisse di fame. I kristang dovevano avere qualche
rifornimento di cibo ruhar, per i prigionieri che catturavano. «Come sta
l'altro soldato, è tutto intero?» Non aveva un bell'aspetto l'ultima
volta che l'avevo visto dentro quel mucchio di mattoni.

«Lei sta bene, sì. Era ferita, ma si è ripresa del tutto. Il suo
velivolo d'assalto è stato danneggiato in orbita, ma è stato salvato da
una delle nostre astronavi.»

Non avevo idea che l'altro soldato fosse una femmina. Con il giubbotto
antiproiettile, il casco e la visiera, era impossibile capirlo. Ed era
mezza sepolta dai detriti del muro che avevamo fatto saltare in aria.
«Per favore le dica che», stavo per aggiungere che mi dispiaceva, cosa
che non era vera perché si trattava di un combattimento, «non era niente
di personale. Non ho avuto il tempo di controllare le sue ferite.»

La Bürgermeister annuì: «Comprensibile, era una situazione di
combattimento. Non ci sono rancori, glielo assicuro. Se mai riuscirò a
contattarla, le riferirò il suo messaggio. Sergente Bishop, l'ho vista
sorpreso di sapere che il soldato ferito fosse una femmina. Avete
femmine tra le vostre fila però».

«Sì, ma è una cosa relativamente recente che le nostre donne siano in
posizioni di combattimento. Anche se le donne hanno prestato servizio
nell'esercito per molti anni e hanno messo le proprie vite a rischio.
Nel combattimento moderno è difficile dire dove siano le prime file.»

«Come hanno reagito i kristang al fatto che gli umani abbiano donne
nelle posizioni di combattimento?»

Era il mio turno di parlare la lingua del corpo, scossi la testa da un
lato all'altro: «Non le dirò niente delle relazioni con i nostri
alleati». Il che era facile, visto che non ne sapevo un accidente.
«Perché lo chiede?»

«Perché i kristang non permettono alle loro donne di prendere parte al
combattimento. O di accedere a una qualsiasi posizione nelle forze
armate. O a qualsiasi ruolo autorevole nella loro società.»

«Non sono affari miei.» Sebbene questo spiegasse i bagni unisex sulle
loro navi: non erano unisex, è che le navi da trasporto truppe non
caricavano mai a bordo kristang di sesso femminile.

«C'è un'espressione umana che ho imparato: ``Conosci il tuo nemico''.
Abbiamo un detto simile tra la mia gente. Le suggerisco anche
d'informarsi sugli esseri che chiamate per errore alleati. I kristang
non ammettono le donne nelle loro forze armate, perché nessuna delle
loro femmine occupa alcuna posizione di autorità in tutta la loro
società. Molto tempo fa, la casta guerriera kristang ha avviato un
programma controllato di riproduzione e ingegneria genetica per ridurre
l'intelligenza delle loro femmine, per renderle più piccole, deboli, più
sottomesse e docili. Il programma d'ingegneria genetica ha anche
alterato il rapporto numerico tra maschi e femmine: in origine la
popolazione era per metà femminile, ora ci sono cinque femmine per ogni
maschio. Per i componenti della casta guerriera kristang, le loro
femmine esistono solo per il piacere, la riproduzione e i lavori
domestici. Quindi, ecco perché le chiedo come reagiscono a vedere le
vostre femmine nel ruolo di soldati. Anche nel ruolo di ufficiali,
piloti e in altre posizioni di prestigio. Posizioni in cui esercitano
autorità sui maschi.»

«Non ho mai incontrato nessun kristang, signora, sono solo un
pivellino.» Non ero certo che ``pivellino'' si traducesse bene in ruhar.
«Sono un soldato di fanteria di basso rango; sono i nostri ufficiali a
gestire i rapporti coi kristang. Non conta quello che i kristang pensano
delle donne nelle nostre forze armate.»

Questa volta, il sorriso della Bürgermeister somigliava più a un
sorrisetto triste. «La pensa ancora così. Come considera i kristang
degli alleati. Forse l'idea dell'allevamento selettivo e della
manipolazione genetica non la preoccupa? Capisco che la storia della
vostra specie include un programma di eugenetica attuato da parte del
governo di un paese chiamato Germania.»

Sapevo che mi stava provocando, ma non riuscii a lasciar correre. «Il
mio paese ha combattuto una guerra contro i nazisti! Contro la Germania,
il paese che chiamiamo Germania», aggiunsi, non sapendo quanto ne
capisse di storia umana. «È stato molto tempo fa. Quello che ha detto
sembra orribile, se è vero.» Pensai a quanto ero stato incazzato in
Nigeria, quando un gruppo di fanatici, perdenti, selvaggi e ignoranti
aveva attaccato una scuola femminile perché pensavano che fosse un
peccato per le ragazze avere un'istruzione. Vedere i corpi bruciati di
quelle ragazzine mi aveva riempito di rabbia, avrei voluto fare
irruzione nella giungla e uccidere tutti i vigliacchi che avevano usato
armi e bombe contro delle bambine indifese e i loro insegnanti. Se solo
fossi riuscito a trovarli. Che era il vero problema mentre eravamo lì.

Cazzo, avevo odiato la Nigeria.

\hfill\break

Nelle quattro settimane successive, la Bürgermeister mi diede un sacco
di informazioni: ben poche buone notizie per gli umani, se diceva la
verità. Ci vorrebbe troppo tempo per ripetere quello che disse, o quello
che dissi io al maggiore Perkins, quindi lo riassumerò. Alcune cose mi
lasciarono con un senso di nausea, e non ero l'unico. Vidi la faccia del
maggiore Perkins sbiancare quando gliene raccontai un po'.

Esporre la Terra all'invasione dei kristang non era stato l'unico
effetto, o l'unico importante, dello spostamento del wormhole. Un intero
quadrante della galassia ne aveva risentito, gettando la società
kristang nel caos. Caos e guerra civile. La casta guerriera kristang non
era una singola entità, era composta da clan che erano in corsa per il
potere e combattevano tra loro. Quando il wormhole vicino alla Terra si
era attivato all'improvviso, ci era voluto un po' di tempo prima che i
thuranin si mettessero a esplorarlo, avendo cose più importanti da fare.
Prima di tutto avevano inviato una sonda robotica nella loro area, per
assicurarsi che l'altra estremità del wormhole non fosse lontana dalla
galassia della Via Lattea, o all'interno di una stella, cosa che era
noto potesse avvenire. Dopo avere stabilito che il wormhole era sicuro
da usare, i thuranin ci avevano inviato un paio di astronavi, poi, di
fatto, avevano deciso con un'alzata di spalle che dalle nostre parti non
c'era molto che valesse la pena sfruttare. Certo, la Terra si trovava a
una settimana di viaggio, ma era lontana da qualsiasi pianeta ruhar, e
lo spostamento del wormhole aveva presentato opportunità più allettanti
altrove. Sarebbe finita così, per secoli forse, con la Terra che
continuava a vagare per il cosmo beatamente ignorante e solitaria, non
fosse che lo spostamento del wormhole aveva anche colpito in modo serio
il patrimonio, già in degrado, del clan Vento Bianco dei kristang.

Prima dello spostamento del wormhole, il clan Vento Bianco era già in
declino e controllava solo due pianeti nel dominio dei kristang. Dopo lo
spostamento del wormhole, i componenti del clan Vento Bianco non avevano
più effettivo accesso a uno di quei pianeti, ed erano stati presi dalla
disperazione. Alcuni imbecilli della leadership del clan Vento Bianco
avevano deciso che la Terra era la risposta ai loro problemi: se
avessero potuto prenderne il controllo e farne una base utile per
invadere il territorio dei ruhar, avrebbero potuto allearsi con un clan
più forte. Non aveva funzionato, nessun altro clan era interessato alla
Terra come base di partenza e l'attacco rovinoso dei ruhar aveva
danneggiato così tanto le nostre infrastrutture che il clan Vento Bianco
avrebbe dovuto spendere troppo tempo e troppe risorse per ricostruirle,
in modo da far sì che potessimo ricominciare a essere utili per loro.
Perciò, la leadership di Vento Bianco aveva deciso che la cosa migliore
da farsi era affittare soldati umani ad altri clan, ed era per questo
che ci trovavamo su Paradiso. Sempre che si volesse dar credito alla
Bürgermeister.

La Bürgermeister avvertì che, se i kristang avessero seguito il loro
schema consueto, avrebbero presto iniziato a richiedere ai governi della
Terra di applicare misure che favorissero i nostri patroni, a spese
degli umani, se non l'avevano già fatto. Tutto in nome della necessità
di ``servire lo sforzo bellico''. Merda, mi sembrava familiare; era
iniziato ancora prima che lasciassi la Terra, e pensavo che le persone
che protestavano contro gli ordini dei kristang ci fossero andate piano
rispetto allo sforzo bellico. Alla fine, i kristang avrebbero stabilito
quale fosse il governo più oppressivo e violento del pianeta e ne
avrebbero fatto il proprio sicario locale, sostenuto da navi da guerra
kristang, intoccabili in orbita. Era, mi disse la Bürgermeister, una
vecchia formula: si trovano i peggiori perdenti psicopatici antisociali
in una società e si dà loro potere in cambio di una completa,
incondizionata lealtà. Ecco come Hitler e Stalin si erano potuti
spingere così in là. Questa era la Procedura operativa standard dei
kristang.

Poi mi parlò dei patroni dei kristang, i thuranin. Erano una specie di
omini verdi, mi disse, confermando ciò che avevo sentito su Campo Alfa.
Quello che non avevo sentito su Campo Raggi X era che i thuranin erano
cyborg, con computer impiantati nel cranio, e la loro società era
altamente interconnessa. Controllavano le loro navi attraverso i
computer impiantati e di rado parlavano ad alta voce, perché
comunicavano tra loro per lo più tramite un collegamento diretto al
computer. La loro pelle era di un beige verdastro, non verde puro, e
tutti avevano lo stesso incarnato. Una volta non era così, ma la razza
dominante aveva commesso un genocidio nei confronti di tutte le altre,
spazzandole via molto tempo prima. Ora tutti i thuranin erano cloni
della ``razza maestra'', con poche variazioni genetiche rispetto a
quello che la razza maestra considerava ideale. I thuranin non
nascondevano il loro disprezzo nei confronti di qualsiasi specie dotata
di una tecnologia inferiore alla loro e disdegnavano anche i loro stessi
mecenati, i maxolhx. Alla base della coalizione guidata dai maxolhx non
c'erano certo vaghi sentimenti di reciproca fratellanza.

Per quanto arroganti fossero i thuranin in merito alla superiorità della
loro tecnologia, la maggior parte di essa non l'avevano sviluppata con
le proprie forze. L'avevano rubata, che è il modo in cui tutte le specie
avevano salito la scala della tecnologia. I motori di salto che i
kristang usavano a bordo delle loro astronavi erano stati copiati da un
drone jeraptha di cui si erano impossessati. Non erano progettati per
essere usati da un'astronave di grandi dimensioni e le copie kristang
erano di scarsa qualità, motivo per cui le loro astronavi potevano
saltare solo per brevi distanze. Gli stessi thuranin avevano rubato,
trovato o comprato un progetto decisamente migliore per le loro
astronavi, così che potessero viaggiare tra le stelle. Ciascuna specie
aveva rubato o copiato la tecnologia da ogni altra specie, anche le due
al vertice di entrambe le fazioni: i maxolhx e i rindhalu.

A volte i maxolhx e i rindhalu si copiavano a vicenda le risorse
tecnologiche, ma al loro livello la fonte principale di tecnologia
avanzata era di gran lunga la ricerca di macchine lasciate dalle antiche
specie, chiamate Anziani. Gli Anziani erano stati la prima specie
intelligente nella galassia della Via Lattea, e pare che vi avessero
vissuto da soli per milioni di anni, prima di scomparire all'improvviso.
Nessuno sapeva che aspetto avessero o perché se ne fossero andati, né
dove; l'ipotesi era che la loro tecnologia fosse ormai così avanzata che
non avevano più bisogno di un'esistenza fisica. A prescindere dal perché
se ne fossero andati e dove, si erano lasciati dietro dispositivi
incredibili, come la rete dei wormhole. E armi. Strumenti d'inaudita
potenza, che avrebbero potuto essere usati come armi, nelle mani
sbagliate.

I kristang ci avevano detto che i maxolhx si erano ribellati contro i
tentativi dei rindhalu di sopprimere lo sviluppo di specie più giovani,
ma la Bürgermeister mi rivelò che in realtà i maxolhx avevano goduto
dell'educazione dei più vecchi rindhalu, finché non avevano trovato
alcune armi degli Anziani e avevano attaccato i loro maestri. Entrambe
le parti avevano schierato armi degli Anziani, ma la guerra aveva avuto
vita breve. L'uso dei dispositivi degli Anziani aveva risvegliato ``le
Sentinelle'', macchine intelligenti che gli Anziani avevano lasciato per
assicurarsi che nessuno abusasse della loro tecnologia residua o
inquinasse la galassia, o per qualche altro motivo. Le Sentinelle
avevano devastato entrambe le parti, poi erano tornate a dormire, o
comunque da dove erano venute. È chiaro che, anche se avevano
abbandonato la galassia, gli Anziani non volevano che specie inferiori
facessero casino con la loro roba.

L'idea delle Sentinelle in agguato nell'ombra spiegava perché maxolhx e
rindhalu non si stessero scontrando in modo diretto nell'attuale guerra
interstellare: quelle due specie non potevano permettersi di combattersi
a vicenda. Era la stessa ragione per cui gli Stati Uniti e la vecchia
Unione Sovietica non si erano mai scontrati in modo diretto nella Guerra
fredda: mutua distruzione assicurata. Qualsiasi lotta di questo tipo tra
gli Stati Uniti e l'Unione Sovietica, entrambi dotati di migliaia di
testate nucleari, sarebbe potuta degenerare in fretta, finendo fuori
controllo e distruggendo entrambe le parti. Qualsiasi combattimento con
le armi degli Anziani avrebbe risvegliato le Sentinelle. Perciò, i
maxolhx e i rindhalu si servivano delle loro specie clienti nel
combattimento per il territorio e le risorse in tutta la galassia,
sperando di schiacciare l'altra parte in un angolo e nell'irrilevanza.
La guerra era combattuta da noi, specie di soldati semplici, e ogni
specie voleva salire la scala della tecnologia e avere specie clienti
inferiori che combattessero e morissero al proprio posto. In quel
momento, gli umani si trovavano in fondo alla scala. Finché non avessimo
sviluppato o rubato una tecnologia avanzata, saremmo rimasti in fondo e
saremmo stati alla mercé dei kristang.

Nel complesso, se la Bürgermeister mi stava dicendo la verità, l'umanità
era fottuta. Il clan Vento Bianco aveva bisogno di trarre qualche
guadagno dalle risorse che aveva investito nella spedizione sulla Terra,
e la Terra non aveva molto da offrire. Certo, di fatto si trovava nel
territorio dei ruhar, ma era così lontana da ogni altro pianeta da loro
occupato che essi potevano permettersi di ignorare la presenza dei
kristang finché non fossero stati pronti ad affrontarla. Cosa avrebbe
potuto offrire la Terra ai kristang, in particolare al clan Vento
Bianco? Il territorio che gli umani stavano già sfruttando avrebbe
potuto essere utilizzato per sostentare la presenza dei kristang. Se ne
sarebbero potute estrarre materie prime con i mezzi più efficienti
possibili, il che significava non preoccuparsi dei terribili danni
ambientali o degli effetti sull'uomo. Si sarebbero potute affittare
truppe umane per lavori che nessun kristang voleva fare, come presidiare
Paradiso. I Vento Bianco dovevano resistere e rendere la Terra una base
di partenza utile, fino al giorno in cui qualche clan più forte avesse
deciso che il nostro piccolo e schifoso pianeta, all'estremità del
territorio ruhar, valeva qualcosa.

Pensai che forse la Bürgermeister ci stava dando qualche bocconcino
d'informazione utile, tipo riguardo ai wormhole, in modo da indurci a
credere alle sue bugie sui kristang. Non spettava a me decidere a cosa
credere, spettava al quartier generale dell'Unef. Il mio lavoro era
sedermi, bere il tè, ascoltarla e passare le informazioni al maggiore
Perkins. Le credevo allora? Ero riluttante a bermi la sua storia. Disse
che la coalizione rindhalu non interferiva con le specie inferiori dal
punto di vista tecnologico, ed era per questo che la Terra era stata
abbandonata, finché i kristang non avevano costretto i ruhar all'azione.
I maxolhx e i loro clienti erano tutti oppressori malvagi, che
sfruttavano specie inferiori e spogliavano i pianeti di risorse, senza
alcuna preoccupazione per le ripercussioni sulle popolazioni indigene.
Era una bella storia, che faceva sembrare buoni i ruhar e cattivi i
kristang, che è quello che mi aspettavo che il nostro nemico dicesse. A
meno che non stesse dicendo la verità. Merda.

\hfill\break

Tra inutili pattugliamenti e raccolta d'informazioni bevendo tazze di
tè, non c'era molto da fare dalle parti di Teskor, dovevo inventarmi
cose per mantenere la squadra occupata e fuori dai guai. Accanto alla
casa, c'era un piccolo blocco di cemento, o di qualcosa di simile al
cemento, perciò montammo un canestro da basket e giocammo partite due
contro due. Le pattuglie della compagnia passavano abbastanza spesso,
tanto che prendemmo le misure di un intero campo da basket, anche se era
di un metro troppo corto, dipingemmo le linee e costruimmo un secondo
canestro. Il nostro era l'unico campo da basket completo della zona e ci
rese così popolari che le pattuglie pianificavano di fermarsi a Teskor
perché potessimo giocare le nostre partite; a volte due pattuglie si
incontravano al nostro posto di comando per prendersi una pausa,
pranzare e giocare a basket. Potevamo anche giocare a softball,
pallavolo, calcio e ping-pong con l'attrezzatura del pacchetto
ricreativo fornito dalla divisione. Niente tavolo da biliardo, cazzo. E
niente Tv. Certo, non c'erano sport o spettacoli da seguire, ma avremmo
potuto guardarci un film. Niente da fare, i ruhar avevano le Tv e noi
avevamo portato Dvd e altra roba con noi dalla Terra, ma era tutto
inutile, dato che le Tv per criceti e i nostri lettori multimediali non
erano compatibili. C'erano schermi piatti portatili per l'uso a livello
aziendale, le unità al di sotto di quella dimensione dovevano
accontentarsi di tablet e laptop. Non era un gran divertimento
affollarsi in quattro attorno a un iPad per guardare un film; voglio
dire, ci si stufa in fretta.

Non pensate che non facessimo altro che giocare partite e guardare film.
L'Unef aveva introdotto gruppi di fuoco in villaggi isolati per
raccogliere informazioni a terra, per vedere cosa facessero i criceti
locali, e noi pattugliavamo e ci guardavamo attorno e tenevamo occhi e
orecchie aperti. Una mattina, Sanchez infilò la testa nella stanzetta
che usavamo da ufficio, dove stavo scrivendo un rapporto per il plotone:
«Ehi, sergente, c'è qualcosa di strano là fuori».

Salimmo sul tetto, avevamo costruito una scala e un ponte di
osservazione in cima alla casa; ci era sembrata la cosa giusta da fare e
ci aveva tenuti occupati. «Guardi il campo Sud 2», disse passandomi il
binocolo.

Sud 2 era un campo vuoto che era stato mietuto un paio di settimane
prima, avevamo visto le grandi mietitrebbie elettriche rombare,
trebbiare il grano e poi sistemarlo in grandi bidoni. Un convoglio di
camion aveva raccolto il grano tre giorni prima, portandolo fino a una
linea ferroviaria a nord della città. «Ottima osservazione, Sanchez»,
dissi, «questo è strano. L'Unef ci ha mandato qui per vedere se accadono
cose strane. In sella.»

Ci mettemmo tutto l'armamentario di battaglia, salimmo sul nostro
cazzuto Crivee e partimmo a tutto gas lungo la strada. Avremmo potuto
attraversare il paese, la terra lì non era nient'altro che campi
all'apparenza vuoti lasciati a maggese, non fosse che avevo ricevuto
dalla divisione un promemoria che ordinava di farla finita con truppe
annoiate che se ne andavano a spasso fuori dalla strada. C'erano stati
diversi incidenti in cui i Crivee si erano ribaltati, il che non andava
bene, visto che l'evacuazione medica per la Terra non era possibile. E
distruggere terreni agricoli di fronte ai criceti non era un buon modo
per conquistare cuori e menti. Così Sanchez si mantenne sulla strada.
Certo, quando arrivammo a Sud 2, c'erano un paio di grandi macchine
elettriche che rombavano lentamente su e giù per i solchi. Piantando
semi, così sembrava. Lasciammo i nostri Crivee sulla strada, Sanchez e
io percorremmo il campo a piedi in direzione delle macchine. Una di loro
si diresse verso di noi a bassa velocità, mentre il criceto alla guida
ci faceva un amichevole cenno di saluto dall'abitacolo in cima. Sul
retro c'era un dispositivo che faceva buchi nel terreno, lasciava cadere
i semi, poi riempiva il buco. Era strano. Così strano che poi facemmo
visita alla famiglia Cornhut.

Lester stava tenendo banco con quella che sembrava metà del villaggio,
perché il fatto di ricevere con regolarità visite della Bürgermeister
aveva aumentato la sua statura a Teskor. Molte famiglie stavano
partecipando a una grigliata dietro casa sua, c'erano tavoli da picnic,
bambini che correvano in giro, una specie di gioco simile al badminton e
due griglie che si riscaldavano. Lester passò di lato alla casa per
venirmi a salutare, con in mano una spatola e indosso un grembiule
bianco con schizzi di salsa rossa, forse provenienti dal barbecue.
Qualunque cosa stesse grigliando aveva un buon odore, anche se sapevo
che gli umani non potevano mangiarne. «Salve Joe Bishop!», disse con un
genuino sorriso amichevole.

Una parte di me si chiese se non fosse il caso di tornare il giorno dopo
e non interrompere la sua grigliata. Avevo un lavoro da fare, così
dissi: «Lester, abbiamo visto delle machine a campo Sud 2», lui
conosceva le designazioni sulle mappe dell'Unef, «stavano piantando
semi».

«Sì, stiamo piantando», il traduttore incespicò, poi disse: «Equivalente
del grano». «È un problema, Joe Bishop?»

«Quanto impiega quel, ehm, raccolto a crescere, prima della mietitura?»

«Tre punto quattro tre mesi», disse il traduttore. Quella caratteristica
dei traduttori era fastidiosa, erano freddamente matematici. È probabile
che Lester avesse detto qualcosa come ``cinque mesi ruhar'', e il
dispositivo aveva fatto i conti. Non avevo bisogno di una precisione
così rigorosa, cazzo.

«L'evacuazione di questo villaggio è prevista fra due o tre mesi»,
stando all'ultima tabella di marcia dell'Unef, che cambiava di continuo,
«perché state piantando un raccolto ora?» Questo era il genere di
``comportamento strano'' dei criceti di cui la G2 voleva essere
informata. Lester aveva in programma di essere ancora a Teskor dopo la
data di evacuazione. I criceti sapevano qualcosa che all'Unef sfuggiva?

«Perché non c'è motivo di non farlo», spiegò Lester, con un ampio
sorriso a trentadue denti che, date le circostanze, non appariva così
amichevole come intendeva. «I kristang ci permettono di continuare a
spedire prodotti alimentari fuori dai confini del pianeta, ma non
possiamo portare i semi con noi. Così piantiamo. Se lasceremo Teskor
prima di poter mietere questo raccolto, forse potrà farlo qualcun
altro.» Alzò le spalle. «Il mio popolo affronta questa guerra da molto
tempo, Joe Bishop. Abbiamo imparato che nulla è certo. Non è certo che i
kristang terranno questo pianeta, quindi speriamo e pianifichiamo per il
futuro, fino al giorno in cui saremo a bordo di una nave che partirà.»

Aveva senso per me: tenere alto il morale dei criceti. La speranza era
un potente aiuto per il morale. «Lester, qualcosa non mi torna.
Dev'essere costoso spedire cibo tra le stelle. Perché farlo?»

«Ah. Questo richiederebbe una lunga spiegazione, Joe Bishop», si voltò a
guardare gli ospiti della sua grigliata, «ci sono due motivi principali.
Nel nostro mondo natale, che a un certo punto era inquinato come mi
dicono sia il vostro, la popolazione ora è concentrata nelle città,
mentre la maggior parte del pianeta è destinata al terreno erboso e
alberato, ripristinato al suo stato naturale. C'è pochissima
agricoltura, perciò bisogna importare il cibo. Molti dei nostri mondi
primari sono tenuti in questo stato, quindi abbiamo altri mondi
designati per l'agricoltura, come Gehtanu, o per l'industria. I mondi
industriali sono per lo più quei pianeti che avevano vita scarsa quando
siamo arrivati, quindi la nostra industria non devasterà l'ambiente
naturale. L'altro motivo per cui siamo determinati a spedire cibo da
questo pianeta il più a lungo possibile, è che lo spostamento del
wormhole, di cui vi è stato detto, ha impedito agli abitanti del nostro
mondo natale di accedere ad alcuni dei mondi che erano destinati
all'approvvigionamento alimentare. Ora potete capire perché siamo
ansiosi di coltivare più cibo possibile qui, prima di partire?»

«Certo, Lester.» Lo lasciammo alla sua grigliata e io spedii un rapporto
alla catena di comando. La spiegazione di Lester mi era sembrata un
tantino una stronzata, ma l'Unef doveva saperlo comunque. Il nostro
gruppo di fuoco iniziò a fare il giro dei terreni agricoli attorno a
Teskor, spaziando in lungo e in largo per controllare le condizioni dei
campi. Quando parlammo con il plotone, scoprimmo che Teskor non era
l'unico villaggio dove i criceti piantavano raccolti che non sarebbero
stati pronti per la mietitura prima della data programmata per
l'evacuazione di quella zona. Non sapevo se credere a Lester o no,
l'intera faccenda non era sotto la mia giurisdizione. Così feci rapporto
al comando e lasciai che se ne preoccupassero i miei superiori.

\hfill\break

A volte, la Bürgermeister sembrava leggermi nella mente e affrontare
argomenti cui stavo già pensando. Mi fece sospettare che i ruhar
stessero in qualche modo origliando quello che dicevamo nel nostro posto
di comando, anche se si suppone che il plotone l'avesse bonificato dalle
cimici. Quel giorno, ero di umore cupo, in sintonia col tempo
atmosferico. Aveva piovuto per cinque giorni e i satelliti meteorologici
prevedevano altri due giorni di pioggia forte, seguiti da una settimana
di rovesci intermittenti. Meraviglioso. I ruhar avevano un nome per il
tempo come questo, lo chiamavano schlumpernur, o almeno così suonava,
solo più stridente. La traduzione era ``coperta umida e ammuffita'', e
noi tutti la trovammo così appropriata che la parola entrò in fretta
nello slang militare degli Stati Uniti. Lo schlumpernur aveva fatto sì
che il villaggio cancellasse la programmazione di un festival del
raccolto, dal momento che nessuno si sentiva, sapete, ``festivo'', e
tenere solo un ``-al'' sarebbe stato triste. I componenti del mio gruppo
di fuoco erano di malumore per essere stati rinchiusi al comando, specie
perché i soldati non possono stare davvero rinchiusi quando il tempo è
brutto, infatti dovevamo comunque uscire di pattuglia e controllare le
condizioni dei campi e inviare rapporti al quartier generale del plotone
e dare l'impressione che stessimo facendo qualcosa di utile. Avevamo
guardato tutti i video a disposizione almeno due volte e, con il tempo
di merda, le pattuglie che passavano da Teskor non si fermavano per una
partita a basket, softball o pallavolo. Cercavamo di attirarle
proponendo una partita a freccette, ma tutti avevano un bersaglio e
nessuno di quelli che passavano voleva lasciare il suo Crivee caldo e
asciutto. Quindi, noi quattro eravamo bloccati insieme ad ascoltare una
dozzina di volte le reciproche storie. Questo era il lato negativo
dell'essere fuori in autonomia; certo, non avevamo un ufficiale che ci
teneva il fiato sul collo, ma ci sentivamo terribilmente soli. Gli
zPhone aiutavano, parlavamo e chattavamo con chiunque sul pianeta. Io
facevo un fischio a Cornpone una volta al giorno per vedere cosa stesse
facendo: era stato assegnato a uno dei plotoni di reazione rapida, il
che sembrava eccitante, ma lui mi disse che non facevano altro che
esercitarsi, esercitarsi e ancora esercitarsi un altro po'; i criceti
non stavano facendo nulla di ostile, quindi non c'era alcunché a cui
reagire in modo rapido. Nel complesso, sembrava quasi annoiato come me.
Concordammo sul fatto che tutti i nostri videogiochi sulla guerra
interstellare ci avevano mentito. E mandai una breve nota a Shauna. Lei
mi mandò una breve nota a sua volta, che presi come un indizio del fatto
che era occupata e non c'era motivo di chattare oltre, a meno che non ci
fossimo visti di persona. Improbabile, dato che era di stanza a
milleseicento chilometri da noi. Perciò, ero di umore cupo quando la
Bürgermeister versò l'acqua calda per il tè. «Joe Bishop, sembra
infelice», osservò.

Feci un'alzata di spalle. «Questo schlumpernur ha abbattuto tutti»,
dissi con l'accenno di un sorriso, perché le scintillarono gli occhi
quando dissi schlumpernur. «E non è che lei mi dia le più felici delle
notizie quando c'incontriamo.» Supponendo che stesse dicendo la verità,
cosa che il quartier generale dell'Unef riteneva fosse scontata.

«Mi ha detto che ha una famiglia sulla Terra?»

Non c'era niente di male nel dirle quello che già sapeva, così annuii:
«I miei genitori e mia sorella».

«Si sentirebbe meglio se avesse loro notizie, no?»

«Certo.» In Nigeria, potevamo usare Skype e ogni tanto anche
video-chattare, almeno un paio di volte a settimana. E con internet
potevamo tutti seguire le notizie a casa, sentirci connessi. «Ho scritto
delle lettere», registrato messaggi video per l'esattezza, «ma finora,
non abbiamo ricevuto risposte.» Inviavamo i messaggi al quartier
generale dell'Unef perché fossero compressi e trasmessi ai kristang, ma,
sebbene le spedizioni di cibo e altri rifornimenti arrivassero con
regolarità in orbita, fino ad allora non avevo ricevuto messaggi da
casa, né notizie di alcun tipo. Era strano. E preoccupante.

«Non ne riceverete», disse lei, guardandomi con attenzione al di sopra
della sua tazza di tè. «I kristang non trasmettono i vostri messaggi a
casa. E non vi trasmetteranno alcun messaggio qui. Non vogliono che
sappiate che cosa sta accadendo sul vostro pianeta. I vostri capi qui lo
sanno di sicuro.»

Merda. Non appena mi diede quelle deprimenti informazioni, la
pioggerella costante si trasformò in un altro acquazzone. La mia
giornata andava di bene in meglio. O aveva deciso che avessi ricevuto
abbastanza brutte notizie per quel giorno, o il tempo aveva depresso
anche lei, perché la Bürgermeister cambiò argomento e, nell'ora che
seguì, parlammo della nostra infanzia. Non ci vidi nulla di male nel
raccontarle noiose storie su com'ero cresciuto nel Maine rurale. La vita
della società ruhar non suonava così diversa dalla vita sulla Terra, non
fosse che avevano una tecnologia incredibile, erano sparpagliati su un
numero indefinito di pianeti, e la sua specie era in guerra da prima che
lei, i suoi genitori, o i suoi bis-bis-bis-bisnonni nascessero,
probabilmente da molto prima. Pensarci non migliorò il mio umore.

\hfill\break

Chen doveva avermi letto nella mente perché, mentre stavamo cenando
quella sera -- pollo alla King per inciso -- mi chiese: «Sergente, ha
idea di quando riceveremo messaggi da casa?».

«No», abbassai gli occhi sul piatto per nascondere uno sguardo colpevole
sul mio viso, «ne so quanto te.»

«Merda. Una nuova nave da trasporto è arrivata dalla Terra ieri, è
ovunque in rete, speravo che la nave portasse lettere, o almeno un po'
di notizie.»

«Sarebbe bello avere notizie», concordò Sanchez, «sono preoccupato per
la mia gente.»

«Tutti lo siamo», ammisi.

«Sapete qual è la cosa peggiore?», chiese Baker. «Di questo dislocamento
non si vede la fine. Quando sono andato in Nigeria, sapevo che sarebbe
durata dodici mesi e che poi sarei tornato a casa. Anche durante la
seconda guerra mondiale, i ragazzi sapevano che c'era una fine in vista:
vinci la guerra e te ne puoi andare a casa. Avrebbero potuto volerci
anni, ma c'era una fine. Anche se le cose si mettevano davvero male, i
ragazzi sapevano che, se avessero resistito e fossero sopravvissuti,
prima o poi sarebbero tornati a casa. Questa guerra dura da migliaia di
anni. Non c'è fine. Non c'è strategia per la vittoria, non ha senso dire
che la missione è compiuta ed è ora di tornare a casa. Ora che ci penso,
nel messaggio in cui mi dicevano che ero stato dislocato fuori dal
pianeta, non specificavano quanto sarebbe durato il dislocamento.»

«Cacciamo i criceti dal pianeta e ce ne andiamo a casa, giusto?», chiese
Chen speranzoso.

«Forse.» Sanchez spinse il suo pollo alla King in giro per il piatto. «A
meno che le lucertole, intendo», mi lanciò uno sguardo preoccupato, «i
kristang, non abbiano bisogno di noi da qualche altra parte. Siamo già
in ballo, giusto? E siamo addestrati, e prima di allora avremo fatto
esperienza. Per i kristang potrebbe avere molto più senso ridislocare
noi che rispedirci in astronave fino a casa e prendere nuove unità.»

«Merda», Chen riassunse i sentimenti di tutti sull'argomento. «Siamo
appena arrivati qui e ho già voglia di andarmene a casa.»

Dovevo mettere fine a quei discorsi cupi. «Ragazzi, guardate. Piacerebbe
anche a me sapere quando torneremo a casa. Spuntare i giorni dal
calendario finché non ci imbarchiamo. Il piano di evacuazione qui dura
tredici mesi, immaginiamo un altro paio di mesi in più: è probabile che
dovremo ricostruire le infrastrutture o qualcosa del genere. Tenerci qui
è dispendioso per i kristang, devono spedirci tutto il cibo per più di
un migliaio di anni luce, giusto? Dopo che avremo compiuto questa
missione, penso che ci porteranno a casa. Andiamo, in questa guerra,
quanti incarichi possono avere per noi i kristang? Anche con i nostri
nuovi giocattoli, non siamo qualificati per un vero combattimento.»

«Sì, ok, credo di sì», Sanchez annuì poco convinto.

«Questo cazzo di schlumpernur ha depresso tutti.» Guardai la pioggerella
fuori dalla finestra. «Se splendesse il sole e ricevessimo qualche
notizia da casa, staremmo tutti alla grande. Siamo fuori in autonomia,
senza ufficiali che ci tengono il fiato sul collo, qui è tutto a posto.
Chen, prendi la scatola della torta dal frigo», quella prelibatezza era
stata inviata dal plotone tre giorni prima, «e festeggiamo. Il meteo
annuncia che questa pioggia finirà dopodomani.»

\hfill\break

Lo schlumpernur era sfumato in piovaschi sparsi prima che il maggiore
Perkins si presentasse per ricevere informazioni fresche da me la
mattina successiva. Ci incontrammo nel suo Crivee, in modo che il mio
gruppo di fuoco potesse restare al riparo mentre pioveva. «Signora, non
so se la maggior parte delle cose che mi sta dicendo siano stronzate o
no, ma non sta mentendo sul fatto che non abbiamo notizie da casa. Il
quartier generale dell'Unef sta ricevendo comunicazioni dalla Terra?»

L'espressione addolorata sul viso del maggiore mi disse tutto quello che
avevo bisogno di sapere.

«Questo è fuori dalla tua giurisdizione, sergente, e dalla mia.»

«I miei ragazzi», accennai alla casa con la testa, «hanno fatto domande
in proposito.»

Mi guardò con fare aggressivo: «Non avrai mica detto qualcosa?».

«No, signora. Sanno che mi sto incontrando con la Bürgermeister e sanno
che lei fa parte della divisione dell'intelligence, e sanno fare due più
due, ma non ho detto loro niente, e loro non l'hanno chiesto. Non hanno
bisogno che sia io a dire che non sappiamo un bel niente dalla Terra. È
in tuta la rete zPhone.» Con gli zPhone, voci che sarebbero rimaste
confinate a una sola unità si erano sparse per tutto il pianeta. Nessuna
delle persone con cui avevo parlato aveva ricevuto un solo messaggio
dalla Terra. E Cornpone mi aveva detto che un tizio dei rifornimenti
sapeva che le spedizioni di cibo dal nostro pianeta erano state scarse
di recente. E tra le provviste che arrivavano allora c'erano semi, come
se i kristang si aspettassero che ci coltivassimo da noi il cibo.

«Sergente», il maggiore distolse lo sguardo e lo rivolse alla pioggia,
«ne so quanto te. Il quartier generale dell'Unef non mi dice più di
quanto abbia bisogno di sapere. Fai attenzione a quello di cui parli in
rete», indicò il tetto del Crivee, «i nostri amici sono in ascolto.»

\hfill\break

Alle quattro del mattino, un paio di giorni dopo, il mio zPhone squillò.
Avere uno zPhone era grandioso sotto molti punti di vista, ma sotto
altri non lo era affatto: il comando poteva raggiungerti in qualunque
momento, ovunque. Raggiungere te personalmente, non chiamare un
operatore radio che doveva poi venirti a cercare. Feci dondolare le
gambe sul pavimento e mi sedetti dritto, qualcuno mi aveva detto che ti
fa sembrare più vigile stare dritto quando ti svegliano di brutto da un
sonno profondo. «Qui sergente Bishop.»

«Bishop, qui è il tenente Charles. Il tuo gruppo di fuoco leverà le
tende oggi per venire qui, bisogna che abbiate sbaraccato e siate in
movimento per le 9:00.»

Questo mi svegliò del tutto. «Cosa succede, signore?»

«Avrete istruzioni quando arriverete qui.»

«Sì, signore», risposi, ma aveva già riagganciato.

I componenti del mio gruppo di fuoco ebbero la stessa reazione alla
``checazzoè'' che avevo avuto io: ne sapevano quanto me, cioè niente. In
rete non girava voce che ci fossero problemi, controllammo i siti
dell'Unef e le notizie erano del tutto normali, considerando che eravamo
su un pianeta alieno. Non ci volle molto, prendemmo le armi, ma
lasciammo tutto ciò di cui avrebbe avuto bisogno un futuro gruppo di
fuoco per occupare il posto, compreso il tavolo da ping-pong che avevamo
costruito. Dopo avere sguazzato per l'ultima volta nelle pozzanghere di
fango lungo la strada principale, svoltammo e lasciammo Teskor prima che
la maggior parte dei criceti fosse uscita dal letto. Mi rammaricai di
non avere avuto l'opportunità di salutare la famiglia Cornhut.

\hfill\break

Sentii puzza di bruciato quando arrivammo alla base del plotone, non
appena la vidi: «Maggiore Perkins? Ci hanno appena fatto ritirare da
Teskor, signora» dissi, pur immaginando che lo sapesse già. Perché con
tutta probabilità l'aveva disposto lei. Non mi sono mai fidato dei tipi
dell'intelligence.

«Bishop, facciamo due chiacchiere.» M'indicò con un cenno il campo
d'aviazione. Una volta che fummo fuori portata d'orecchio, spiegò: «Sono
stata io a ordinare che vi ritiraste da Teskor. I kristang sanno che
qualcuno ci ha passato informazioni e sono incazzati per questo, e
stanno infilando il naso in giro e si stanno avvicinando a te. Devi
nasconderti per un po', così ti trasferiscono. Solo tu, non la tua
squadra».

«Dove, signora?»

«Il Vettore di carico. Ti dirò la verità, non abbiamo un incarico per te
in questo momento, a parte sparire e tenere il tuo nome fuori dai canali
di comunicazione. Se ti fa sentire meglio, mi è stato detto lo stesso,
mi stanno riassegnando a una base logistica nel settore indiano, come
collegamento.» Sbuffò. «Non parlo una sola parola di hindi.»

«Cazzo. Tutto questo perché abbiamo ascoltato i criceti? I kristang
pensavano che non avremmo parlato con loro per tutto il tempo che
saremmo rimasti qui? Devono sapere che i criceti approfitterebbero di
ogni occasione per far fare brutta figura alle lucer\ldots{} ai
kristang, per metterci i bastoni tra le ruote.»

«Non è soltanto il fatto di avere parlato con loro, i kristang si
aspettano che in qualunque esercito le voci girino. Il problema è che
sanno che il comando dell'Unef ha preso sul serio i criceti. Tu non sei
la nostra sola fonte diretta d'informazioni, ma posso dirti che la tua
Bürgermeister è considerata la fonte più preziosa, e il comando
dell'Unef si sta cagando addosso, perché pensa che la maggior parte
delle cose che ci ha detto siano vere.» Mi lanciò uno sguardo diretto:
«Non devi ancora dire una parola di tutto questo. È top secret».

«Sì, signora.»

Fece una lunga espirazione, si accarezzò la tasca della camicia, poi
scosse la mano con disgusto: «Essere fuori nella galassia alla fine mi
ha fatto smettere di fumare, quando la mia scorta è finita ho smesso di
punto in bianco. I kristang considerano il tabacco un bene di lusso che
non vogliono trasportare attraverso gli anni luce. A volte avrei voglia
di sparare a qualcosa».

«Eisenhower smise di punto in bianco, signora, dopo il primo attacco di
cuore.» Questa non mi uscì esattamente come l'avevo intesa.

«Se Ike lo fece, mi sta bene. E lui aveva scelta, giusto? Io no. Forse
la nicotina avrebbe dovuto essere nella lista delle medicine
essenziali», disse il maggiore con una smorfia.

Avrei voluto dimostrare maggior interesse, ma non ci riuscivo. Provavo
empatia, ma, poiché non avevo mai dovuto perdere l'abitudine di fumare,
bere o drogarmi, non potevo davvero comprendere cosa stesse passando,
capite che intendo? Avevo scelto la via facile di non andare in fissa
con niente a priori. Non sapendo cosa dire, feci un amichevole suono
``uhm'' e aspettai che fosse lei a parlare.

«La medicina è ciò che ci ha messo in guai seri comunque. Hai ragione,
alle lucertole non frega un cazzo delle voci che girano. La cosa che li
ha fatti incazzare è che abbiamo accettato l'offerta dei ruhar di
fornirci assistenza medica avanzata. Hai sentito dell'incidente della
Poiana il mese scorso? Due morti, quattro feriti, uno dei feriti ha
perso una gamba, un altro si è rotto la spina dorsale. E i kristang non
volevano riportarli con un volo di evacuazione medica sulla Terra. I
ruhar hanno saputo dei feriti, sono rimasti sorpresi di sentire che gli
umani non hanno la capacità di far ricrescere gli arti o il tessuto
nervoso, loro utilizzano quella tecnologia da così tanto tempo che la
danno per scontata. Come esperimento, abbiamo mandato alcuni casi di
feriti gravi in un ospedale ruhar e sono stati curati, e tra loro
c'erano anche due dei feriti nell'incidente della Poiana. I medici ruhar
ci hanno messo un po' a adattarsi alla nostra diversa biochimica, ma la
loro è basata sul Dna come la nostra. I rapporti dicono che alla fine ci
si aspetta un completo recupero di tutti i feriti, anche di quello che
ha perso una gamba, destinata a ricrescere. I kristang si sono infuriati
quando hanno saputo che avevamo accettato le cure mediche dei ruhar,
hanno rampognato il comando Unef per aver fraternizzato con il nemico.
Ho sentito che la conversazione si è scaldata: il generale Meers ha
detto loro che non saremmo stati costretti a mandare la nostra gente
negli ospedali ruhar, se solo le lucertole avessero condiviso la loro
tecnologia medica.» Sorrise. «Il vecchio non arretra di fronte a
nessuno, umano o lucertola che sia.»

«Sono nei guai, signora?»

«Cazzo, no, Bishop. Mantieni un profilo basso finché questa cosa non si
sgonfia da sé nel giro di un paio di settimane, mesi forse. Pensala come
un'opportunità di vedere qualcosa di più del pianeta.»

Non ci credetti neanche per un secondo.

\hfill\break

Dopo avere salutato il mio gruppo di fuoco, presi una serie di voli per
il Vettore: il viaggio fu lungo e noioso, ed ero di pessimo umore. Non
avevo fatto niente di male, ma, se l'Unef avesse avuto bisogno di un
capro espiatorio per i kristang, si sarebbero aspettati che mi prendessi
tutta la colpa. Una riduzione di grado sembrava probabile, non ero stato
un sergente a lungo e mi ero appena abituato, mi sarebbe mancata la
sensazione di essere a capo di una squadra. Avevamo fatto un buon lavoro
a Teskor, ci eravamo tenuti fuori dai guai, avevamo stabilito buoni
rapporti con i nativi, avevamo raccolto informazioni solide e le avevamo
trasmesse ai superiori. E avevamo costruito un campo da basket che aveva
fatto del nostro posto di comando un luogo popolare, dove le pattuglie
potevano prendersi una sosta, cosa che faceva bene al morale. Non c'era
stato alcun serio attrito all'interno del mio gruppo di fuoco, non più
di quanto ci si aspetterebbe da quattro ragazzi bloccati insieme in
mezzo al nulla con poco da fare. Ripercorrendo il tempo trascorso a
Teskor, mi era sembrato quasi idilliaco, tranquillo. Eravamo atterrati
su un pianeta alieno, nemico, avevamo stabilito il controllo e stavamo
facendo la nostra parte nella guerra, mostrando ai kristang che potevamo
dare il nostro contributo, per quanto piccolo potesse essere. Sentivamo
di avere un senso di scopo, di realizzazione, tutti aspettavamo con
ansia il giorno in cui gli abitanti di Teskor si sarebbero imbarcati sui
mezzi di trasporto e avrebbero preso parte all'evacuazione. Sarebbe
stata una missione compiuta per la nostra squadra integrata di
osservazione.

Ripercorrendo il tempo da noi trascorso a Teskor dopo avere iniziato a
ricevere informazioni inquietanti dalla Bürgermeister, la visione che ne
avevo era diventata meno idilliaca. Apprendere che l'umanità era stata
raggirata e indotta a istituire una Expeditionary Force, che l'Unef non
era altro che un soldatino dei kristang, che le lucertole consideravano
la Terra un trofeo di guerra da sfruttare come volevano uccise il mio
entusiasmo per ogni sensazione di successo. Il fatto di non poter
condividere informazioni con il mio gruppo di fuoco aveva interposto una
distanza tra noi, loro sapevano che stavo incontrando un funzionario
ruhar di alto rango, avevano visto un ufficiale dei servizi segreti
dell'Unef in visita il giorno dopo che avevo incontrato i ruhar, e
sapevano fare due più due. Inoltre, va da sé, giravano voci. Alcune
riguardavano informazioni che avevo ricevuto dalla Bürgermeister, e mi
davano da pensare riguardo alla sicurezza informativa del quartier
generale dell'Unef. Mi faceva incazzare il fatto di dover trattare le
informazioni come top secret, quando giravano già voci in merito. Non
potevo nemmeno dire ai miei ragazzi quali voci fossero una totale
stronzata e quali contenessero almeno un fondo di verità. Di totali
stronzate ne giravano molte.

Ora i componenti del mio gruppo di fuoco avevano un nuovo sergente ed
erano di stanza al quartier generale del battaglione, chiedendosi se
avessero fatto qualcosa di sbagliato. Nessuno credeva alla storia del
cazzo che la nostra squadra integrata di osservazione fosse stata
ritirata da Teskor perché la missione era finita: quegli stupidi del
quartier generale dell'Unef non erano stati abbastanza intelligenti da
fare il cambio con un'altra squadra integrata di osservazione, avevano
semplicemente abbandonato Teskor. Certo, non appariva per niente
sospetto, le squadre integrate di osservazione erano ancora attive in
tutto il settore, tranne che nell'unico villaggio in cui avevo
incontrato la Bürgermeister. Ecco quanto era stata sottile la copertura
dell'Unef: il caposquadra del Dumbo che portai al Vettore, un tizio che
non avevo mai incontrato in vita mia, mi chiese cosa avessi fatto per
costringere l'Unef a ritirare la nostra squadra da Teskor. Giravano voci
di uno scandalo, voci così succulente da inondare il pianeta alla
velocità della luce. Tutte le misure di InfoSec applicate dall'Unef
potevano essere neutralizzate da due tizi che condividono un
pettegolezzo sui loro zPhone.